\Needspace{6\baselineskip}
\begin{center}\SourceLeaf{248}\pdfbookmark[2]{176.}{article-176}176.\end{center}
Just as, for any fixed values of the unknowns, each system of values of the functions $V$, $V'$, $V''$ etc. has a definite probability before the observation is made, so conversely, once definite values of the functions have been obtained from the observations, a definite probability will attach to each system of values of the unknowns from which they could have arisen: clearly, the systems for which the event that did occur was more to be expected must be regarded as the more probable. The estimation of this probability rests on the following theorem:

\emph{If, under a hypothesis $H$, the probability of a given event $E$ is $=h$, while under another hypothesis $H'$, which excludes the first and is in itself equally probable, the probability of the same event is $=h'$: then I say that, when the event $E$ has actually occurred, the probability that $H$ was the true hypothesis will be to the probability that $H'$ was the true hypothesis as $h$ is to $h'$}\HistoricalNoteEight{Manuscript note:] If the hypotheses $H$, $H'$ had in themselves (i.\,e. before $E$ occurred or before it became known that it had occurred), unequal probabilities $\mu$, $\mu'$, then we would have to assign them, after the occurrence of $E$, probabilities proportional to the products $\mu h$, $\mu' h'$.}.

To prove this, suppose that by distinguishing all the circumstances on which it depends whether $H$ or $H'$ or some other hypothesis holds, and whether the event $E$ or another event must occur, we obtain a system of distinct cases, each of which in itself (i.\,e. as long as it is uncertain whether the event $E$ or another event will occur) is to be regarded as equally probable, and that these cases are distributed as follows,
\begin{center}
\small
\begin{tabular}{c|c|c}
\shortstack{so that among them\\there are} & \shortstack{in which the hypothesis\\must hold} & \shortstack{under such conditions\\that the event must be} \\
\hline
$m$ & $H$ & $E$ \\
$n$ & $H$ & different from $E$ \\
$m'$ & $H'$ & $E$ \\
$n'$ & $H'$ & different from $E$ \\
$m''$ & different from $H$ and $H'$ & $E$ \\
$n''$ & different from $H$ and $H'$ & different from $E$
\end{tabular}
\end{center}
Then we have $h=\dfrac{m}{m+n}\mathpunct{\raisebox{.45ex}{,}}$ $h'=\dfrac{m'}{m'+n'}\mathpunct{\raisebox{.45ex}{;}}$ further, before the event was known, the probability of \SourceLeaf{249}hypothesis $H$ was $=\dfrac{m+n}{m+n+m'+n'+m''+n''}\mathpunct{\raisebox{.45ex}{,}}$ but after the event is known, when the cases $n$, $n'$, $n''$ cease to be possible, the probability of the same hypothesis will be $=\dfrac{m}{m+m'+m''}\mathpunct{\raisebox{.45ex}{;}}$ similarly the probability of hypothesis $H'$ before and after the event is expressed, respectively, by $\dfrac{m'+n'}{m+n+m'+n'+m''+n''}$ and $\dfrac{m'}{m+m'+m''}\mathpunct{:}$ since, therefore, the hypotheses $H$ and $H'$ are assumed equally probable before the event is known, we will have $m+n=m'+n'$, from which the truth of the theorem follows immediately.

Now insofar as we assume that, apart from the observations $V=M$, $V'=M'$, $V''=M''$ etc., no other data for determining the unknowns are available, and thus that all systems of values of these unknowns were equally probable before those observations, the probability of any particular system after those observations is clearly proportional to $\Omega$ itself. This is to be understood in the sense that the probability that the values of the unknowns lie respectively between the infinitely close limits $p$ and $p+\mathrm{d}p$, $q$ and $q+\mathrm{d}q$, $r$ and $r+\mathrm{d}r$, $s$ and $s+\mathrm{d}s$ etc., is expressed by $\lambda\Omega\,\mathrm{d}p\,\mathrm{d}q\,\mathrm{d}r\,\mathrm{d}s$ etc., where $\lambda$ is a constant quantity independent of $p$, $q$, $r$, $s$ etc. And clearly $\dfrac{1}{\lambda}$ will be the value of the integral of order $\nu^{ti}$: $\int^{\nu}\Omega\,\mathrm{d}p\,\mathrm{d}q\,\mathrm{d}r\,\mathrm{d}s\, .\, .\, .\, .$, each of the variables $p$, $q$, $r$, $s$ etc. ranging from $-\infty$ to $+\infty$.
